\chapter{Devices and Sensors}
\label{ch:devices-sensors}

This chapter covers two namespaces that share a directory prefix but almost nothing else:
\cnans{Microsoft::Devices::Sensors} is a real, historical Windows Phone~7 API surface CNA
faithfully ports; \cnans{CNA::Devices} is an entirely CNA-original, \texttt{NOXNA} desktop
utility layer with no XNA or WP7 equivalent at all.

\section{Microsoft::Devices::Sensors: a real WP7 API, engineered for real concurrency}

\cnaclass{SensorBase<TSensorReading>} is the shared template base for every sensor in this
namespace (\cnaclass{Accelerometer}, \cnaclass{Gyroscope}, \cnaclass{Compass},
\cnaclass{Motion}), providing \texttt{CurrentValue}, \texttt{IsDataValid},
\texttt{TimeBetweenUpdates} (defaulting to 2 milliseconds, matching the real WP7 API), the
public \texttt{CurrentValueChanged} event, and abstract \texttt{Start()}/\texttt{Stop()}.
What makes this namespace worth a dedicated section in this book is not its API surface,
which is small, but the genuine concurrency-safety engineering behind it --- because, unlike
most of CNA's XNA-facing surface, a sensor's \texttt{CurrentValueChanged} callback can
legitimately fire from whatever thread SDL's own sensor event-watch mechanism happens to
invoke it on, not necessarily the game's own update thread.

\subsection{A real, found-and-fixed data race, and the design it produced}

An early implementation read \texttt{isDisposed\_} with no lock at all, while the real
SDL-backed sensor classes wrote to it from a background event-watch thread --- a genuine
data race under the C\texttt{++} memory model, not a theoretical one. The fix that grew out
of finding it is more interesting than the fix itself: rather than a single coarse lock, this
namespace now uses a documented, per-instance protocol for safe disposal under concurrent
access --- an atomic \texttt{ClaimDisposalOnce()} so exactly one of two racing
\texttt{Dispose()} callers actually runs cleanup, a \texttt{WaitForDisposalToComplete()} for
the loser to block on, and (added after a later external audit found the first version could
still strand a waiting thread if the winning cleanup itself threw) an RAII
\texttt{DisposalTerminalStateGuard} that unconditionally publishes the disposed state even on
an exceptional exit path. \texttt{Accelerometer} and \texttt{Compass} each additionally track
which specific threads are currently mid-dispatch into a callback, so that a handler which
reentrantly disposes its own sensor instance does not deadlock waiting for itself to finish.

\subsection{A second, subtler bug: an update-throttling integer overflow}

A separate, later external audit found a genuine signed-integer-overflow bug in the update
throttle every sensor uses to honor \texttt{TimeBetweenUpdates}: comparing an elapsed
duration against the configured interval by promoting both to a common, finer time
resolution multiplies the interval's own tick count by 100 --- which overflows a signed
64-bit integer outright when \texttt{TimeBetweenUpdates} is set to \texttt{TimeSpan.MaxValue},
confirmed directly under UndefinedBehaviorSanitizer. The fix casts the \emph{elapsed}
duration down to the coarser resolution instead of promoting the interval up --- a direction
that can never overflow, since a signed 64-bit count of 100-nanosecond ticks can represent
roughly 29{,}000 years, far beyond any realistic elapsed wall-clock time a running program
would ever measure.

\subsection{Per-sensor reality: what is genuinely backed by hardware}

\cnaclass{Accelerometer} is the most complete of the four: genuinely SDL3-backed
(\texttt{SDL\_SENSOR\_ACCEL}) on desktop, Android, and iOS alike --- a deliberate design
decision, not a Windows Phone accident, since desktop accelerometer hardware (a laptop's own
motion sensor, for instance) is treated as a first-class supported case rather than a
mobile-only special case. It also carries the one genuinely real \texttt{State} property in
this namespace, and the legacy WP7~7.0 \texttt{ReadingChanged} event --- deprecated since WP7
7.1 in favor of \texttt{CurrentValueChanged}, but still present and still raised immediately
after \texttt{CurrentValueChanged} for the same reading, exactly matching real WP7's own
still-shipping deprecated-but-functional behavior. \cnaclass{Gyroscope} follows the identical
shape and the identical throttling fixes. \cnaclass{Compass}, by contrast, has no SDL3
backing at all --- SDL3 exposes no magnetometer API on any platform --- so it is only
genuinely real on Android, through a dedicated native backend built directly on the NDK's own
rotation-vector and magnetic-field sensors; on every other platform, \texttt{IsSupported}
correctly reports false and \texttt{Start()} fails with \cnaclass{SensorFailedException},
exactly as real WP7 hardware without a compass would behave. \cnaclass{Motion} (a fused
attitude/gravity/acceleration reading) shares this same Android-real,
everywhere-else-honest-stub shape.

\cnaclass{VibrateController} is a real WP7 API in its own right: a single, non-static,
process-lifetime shared instance reached only through \texttt{getDefaultProperty()} --- the
private constructor exists specifically because a phone has exactly one vibration motor, and
a second, independently constructed instance could not meaningfully represent anything
different. \texttt{Start(TimeSpan)} throws outside the documented $[0, 5\text{s}]$ WP7 range;
a real \texttt{NOXNA} extension, \texttt{Start(TimeSpan, float intensity)}, exposes SDL3's
graduated rumble-strength parameter directly, since modern hardware supports variable
intensity that a 2010-era WP7 phone's simple on/off buzzer never did. On Android specifically,
a documented, source-verified hardware limitation is worth knowing: SDL3's own Android haptic
backend blends any two independent motor intensities into one blended value before handing it
to the platform vibration service --- true independent left/right rumble on Android is only
reachable through \cnaclass{GamePad}'s own controller-rumble path, never through this class.

\section{CNA::Devices: a NOXNA desktop utility layer}

Distinct from the sensor namespace above, \cnans{CNA::Devices} has no XNA or WP7 precedent at
all --- every class in it is a CNA-original addition, built directly on an SDL3 capability
with no XNA equivalent to preserve fidelity against. \cnaclass{Camera} captures live video
frames from a camera device directly into a \cnaclass{Texture2D}, deliberately poll-based
(matching SDL3's own \texttt{SDL\_AcquireCameraFrame()} shape) rather than forcing a
callback model onto a fundamentally poll-driven API; its first-implementation scope is
explicitly narrow --- a single, first-reported camera device, synchronous permission
polling, RGBA8-only frame delivery. \cnaclass{SystemTray} shows a real system tray icon with
a flat menu, confirmed genuinely desktop-only by reading SDL3's own per-platform tray source
directly (no Android, iOS, or Emscripten backend exists at all --- unlike
\cnaclass{FileDialog}, which turned out to have a real Android backend once actually
investigated, rather than being assumed desktop-only by analogy). \cnaclass{DisplayInfo}
deliberately does not duplicate what \cnaclass{GameWindow} (this book's game-loop chapter)
already exposes --- only the two genuinely new SDL3 capabilities
\cnaclass{GameWindow} has no property for at all: per-window display content scale, and a
window's safe interactive area (relevant on displays with notches or rounded corners).
\cnaclass{MessageBox}, \cnaclass{FileDialog}, \cnaclass{UrlLauncher}, \cnaclass{Locale},
\cnaclass{Clipboard}, \cnaclass{SystemInfo}, and \cnaclass{PowerInfo} round out this layer,
each a thin, real wrapper over the corresponding SDL3 host-integration capability.

\section{Why this namespace matters despite its small XNA surface}

Real XNA's own \cnans{Microsoft.Devices.Sensors} surface is small, and CNA's port of it is
correspondingly modest in scope compared to Graphics or Audio. What makes it worth a chapter
in this book is not size but craft: this is the one namespace in the whole project whose own
documentation most explicitly foregrounds thread-safety as a first-class design concern
(\texttt{docs/devices-thread-safety.md} is described, in the source itself, as the single
consolidated contract every sensor class's own comments defer back to), precisely because it
is the one place in CNA where a callback genuinely can arrive from a thread the game's own
code never spawned.
